
\subsection{Descent direction to escape the saddle point on the antisymmetric branch}
\label{subsec:descent-direction-to-escape}

This subsection uses the transverse teacher imbalance (\cref{assumption:quantitative-unweighted-two-sided-transverse-imbalance}) to construct a
same-density perturbation that escapes from the antipodal saddle.

The imbalance assumption says that the teacher has a coherent first-order
transverse bias across two-sided latitude slices around the antipodal axis. Here we show that this bias is visible to the active odd Hermite features, and
therefore produces a same-density tangent direction that decreases the fitting
loss to second order at the antipodal Dirac point.

\begin{remark}[Comparison with the zero-output saddle argument of \citet{yamamoto2025hessianguided}]
    The algebraic mechanism is the same as the known zero-output saddle calculation in
the matrix-decomposition example (Proposition G.1 in \citet{yamamoto2025hessianguided}): when the
student output vanishes, a perturbation that creates a first-order feature
\(g\) changes the product output only at order \(h^2\), producing the quadratic
cross term
\[
    -\E_z[(h^o(z)^\top g(z))^2].
\]
However, the affine mixture curve used there does not preserve our product
spherical ansatz.  We therefore rebuild the argument
with a Dirac-supported spherical tangent perturbation at the antipodal pair.
The transverse imbalance assumption is precisely the condition ensuring that one such tangent feature is visible to the teacher.
\end{remark}

The first step turns the sliced transverse imbalance assumption into a nonzero active-weighted transverse teacher moment.

\begin{lemma}[Quantitative control of the active-weighted transverse moment]
\label{lem:quantitative-control-weighted-transverse-moment}
\textbf{Assumptions used.}
\cref{def:two-sided-transverse-teacher-moment-measure} and 
\cref{assumption:quantitative-unweighted-two-sided-transverse-imbalance}.

Set \(a=a_\star\), the anchor in that imbalance assumption.  Define
\[
    \kappa_{\mathcal O}(r)
    \coloneq
    \frac12
    \sum_{n\in\mathcal O}
    n\|a_n\|_2^2 r^{n-1},
    \qquad r\in[0,1].
\]
Assume at least one active coefficient \(\|a_n\|_2\) is nonzero.  Since every
\(n\in\mathcal O\) is odd,
\[
    \kappa_{\mathcal O}(r)>0
    \qquad
    \forall r\in(0,1].
\]
Let
\[
    b_a^o
    \coloneq
    \int_{\Sphere^{d_w-1}}
    \kappa_{\mathcal O}(|a^\top w|)
    \Pi_{a^\perp}w
    \,\dd\nu^o(w)
    \in a^\perp .
\]
Equivalently,
\[
    b_a^o
    =
    \int_{[0,1]}
    \kappa_{\mathcal O}(r)\,\dd\mathfrak m_a^o(r).
\]
Define
\[
    \kappa_{\min}(r_{\mathrm{vis}})
    \coloneq
    \inf_{r\in[r_{\mathrm{vis}},1]}
    \kappa_{\mathcal O}(r).
\]
Then
\[
    \kappa_{\min}(r_{\mathrm{vis}})>0
\]
and
\[
    \|b_a^o\|
    \ge
    \langle e_\star,b_a^o\rangle
    \ge
    \kappa_{\min}(r_{\mathrm{vis}})\gamma_\perp.
\]
In particular,
\[
    b_a^o\neq0.
\]
\end{lemma}

\begin{proof}
We first rewrite the weighted transverse moment using the latitude
disintegration.  By
\cref{def:two-sided-transverse-teacher-moment-measure}, for every bounded
Borel function \(\psi:[0,1]\to\mathbb R\),
\[
    \int_{\Sphere^{d_w-1}}
    \psi(r_a(w))\Pi_{a^\perp}w\,\dd\nu^o(w)
    =
    \int_{[0,1]}
    \psi(r)\,\dd\mathfrak m_a^o(r).
\]
Taking
\[
    \psi(r)=\kappa_{\mathcal O}(r)
\]
gives
\[
    b_a^o
    =
    \int_{[0,1]}
    \kappa_{\mathcal O}(r)\,\dd\mathfrak m_a^o(r).
\]
Therefore,
\[
    \langle e_\star,b_a^o\rangle
    =
    \int_{[0,1]}
    \kappa_{\mathcal O}(r)\,
    \dd\mathfrak m_{a,e_\star}^o(r).
\]

Next, since every \(n\in\mathcal O\) is odd, \(n-1\) is even.  Hence
\[
    r^{n-1}>0
    \qquad
    \forall r\in(0,1]
\]
for every active degree \(n\).  Because at least one active coefficient is
nonzero,
\[
    \kappa_{\mathcal O}(r)>0
    \qquad
    \forall r\in(0,1].
\]
Since \(r_{\mathrm{vis}}\in(0,1)\) and \(\kappa_{\mathcal O}\) is continuous,
\[
    \kappa_{\min}(r_{\mathrm{vis}})
    =
    \inf_{r\in[r_{\mathrm{vis}},1]}
    \kappa_{\mathcal O}(r)
    >0.
\]

By
\cref{assumption:quantitative-unweighted-two-sided-transverse-imbalance},
\(\mathfrak m_{a,e_\star}^o\) is a positive finite measure on \([0,1]\).
Thus
\[
    \langle e_\star,b_a^o\rangle
    =
    \int_{[0,1]}
    \kappa_{\mathcal O}(r)\,
    \dd\mathfrak m_{a,e_\star}^o(r)
    \ge
    \int_{[r_{\mathrm{vis}},1]}
    \kappa_{\mathcal O}(r)\,
    \dd\mathfrak m_{a,e_\star}^o(r).
\]
Using the definition of \(\kappa_{\min}(r_{\mathrm{vis}})\), we obtain
\[
    \int_{[r_{\mathrm{vis}},1]}
    \kappa_{\mathcal O}(r)\,
    \dd\mathfrak m_{a,e_\star}^o(r)
    \ge
    \kappa_{\min}(r_{\mathrm{vis}})
    \mathfrak m_{a,e_\star}^o([r_{\mathrm{vis}},1]).
\]
The quantitative strict imbalance gives
\[
    \mathfrak m_{a,e_\star}^o([r_{\mathrm{vis}},1])
    \ge
    \gamma_\perp.
\]
Consequently,
\[
    \langle e_\star,b_a^o\rangle
    \ge
    \kappa_{\min}(r_{\mathrm{vis}})\gamma_\perp.
\]
Finally, since \(\|e_\star\|=1\),
\[
    \|b_a^o\|
    \ge
    \langle e_\star,b_a^o\rangle.
\]
This proves the claim.
\end{proof}

We next encode how a tangent displacement of the antipodal atoms is seen by the
teacher feature.

\begin{definition}[Tangential visibility matrix at an antipodal Dirac pair]
\label{def:tangential-visibility-matrix-dirac-pair}
Let
\[
    \nu_a^\pm
    \coloneq
    \frac12\delta_a+\frac12\delta_{-a}.
\]
For \(s\in a^\perp\), define the Dirac-supported active differential
\[
    \mathcal U_a(s)_n
    \coloneq
    \frac n2
    \sym{a^{\otimes(n-1)}\otimes s},
    \qquad n\in\mathcal O.
\]
Let \(g_{a,s}\) be the corresponding first-order student feature,
\begin{equation}
\label{eq:mf-visibility-first-order-feature}
 g_{a,s}(z):=\sum_{n\in\mathcal O}a_n
       \langle\Hertensor_n(z),\mathcal U_a(s)_n\rangle_{\rm HS}.
\end{equation}
This uses the finite expansion in \eqref{eq:activation-moment-expansion}.
Recall the teacher from \cref{assumption:odd-symmetric-teacher}:
\[
 h^o(z)=h_{\nu^o}(z)
 =\sum_{n\in\mathcal O}a_n\langle\Hertensor_n(z),T_n^o\rangle_{\rm HS},
 \qquad T_n^o=\int w^{\otimes n}\,d\nu^o(w).
\]
Equivalently, define the linear map
\[
    J_a(z):a^\perp\to\mathbb R^{d_h},
    \qquad
    J_a(z)s=g_{a,s}(z).
\]
For the teacher feature \(h^o\), define the positive semidefinite operator
on \(a^\perp\)
\[
    \mathsf B_a^o
    \coloneq
    \E_z
    \left[
        J_a(z)^\top h^o(z)h^o(z)^\top J_a(z)
    \right].
\]
We say that the teacher is tangentially visible at the antipodal pair
\(\{\pm a\}\) if
\[
    \lambda_{\max}(\mathsf B_a^o)>0.
\]
\end{definition}

The nonzero transverse teacher moment implies that some tangent direction has
nonzero correlation with the teacher output.

\begin{lemma}[Active-weighted transverse moment implies tangential visibility]
\label{lem:active-weighted-transverse-moment-implies-visibility}
\textbf{Assumptions used.}
\cref{assumption:on-sphere,assumption:odd-symmetric-teacher},
together with
\cref{def:tangential-visibility-matrix-dirac-pair}.

Let
\[
    b_a^o
    \coloneq
    \int_{\Sphere^{d_w-1}}
    \kappa_{\mathcal O}(|a^\top w|)
    \Pi_{a^\perp}w
    \,\dd\nu^o(w).
\]
Then
\[
    \mathsf B_a^o
    \succeq
    b_a^o(b_a^o)^\top.
\]
Consequently,
\[
    \lambda_{\max}(\mathsf B_a^o)
    \ge
    \|b_a^o\|^2.
\]
In particular, if \(b_a^o\neq0\), then the teacher is tangentially visible at
\(\{\pm a\}\).
\end{lemma}

\begin{proof}
For \(s\in a^\perp\), Hermite orthogonality gives
\[
    \E_z\bigl[g_{a,s}(z)^\top h^o(z)\bigr]
    =
    \sum_{n\in\mathcal O}
    \|a_n\|_2^2
    \left\langle
        \mathcal U_a(s)_n,T_n^o
    \right\rangle_{\HS}.
\]
By the definition of \(\mathcal U_a(s)\),
\[
    \left\langle
        \mathcal U_a(s)_n,T_n^o
    \right\rangle_{\HS}
    =
    \frac n2
    \left\langle
        \sym{a^{\otimes(n-1)}\otimes s},
        T_n^o
    \right\rangle_{\HS}.
\]
Since \(T_n^o\) is symmetric,
\[
    \left\langle
        \sym{a^{\otimes(n-1)}\otimes s},
        T_n^o
    \right\rangle_{\HS}
    =
    \left\langle
        s,
        T_n^o\contr{n-1}a^{\otimes(n-1)}
    \right\rangle.
\]
Using
\[
    T_n^o
    =
    \int_{\Sphere^{d_w-1}} w^{\otimes n}\,\dd\nu^o(w),
\]
and \(s\in a^\perp\), we obtain
\[
    \left\langle
        s,
        T_n^o\contr{n-1}a^{\otimes(n-1)}
    \right\rangle
    =
    \int_{\Sphere^{d_w-1}}
    (a^\top w)^{n-1}
    \langle s,\Pi_{a^\perp}w\rangle
    \,\dd\nu^o(w).
\]
Since \(n\) is odd, \(n-1\) is even, hence
\[
    (a^\top w)^{n-1}
    =
    |a^\top w|^{n-1}.
\]
Therefore
\[
    \E_z\bigl[g_{a,s}(z)^\top h^o(z)\bigr]
    =
    \left\langle s,b_a^o\right\rangle.
\]

Set
\[
    X(z)\coloneq J_a(z)^\top h^o(z)\in a^\perp.
\]
Then
\[
    \mathsf B_a^o=\E_z[X(z)X(z)^\top],
    \qquad
    \E_z[X(z)]=b_a^o.
\]
By Jensen's inequality in the Loewner order,
\[
    \mathsf B_a^o
    =
    \E_z[X(z)X(z)^\top]
    \succeq
    \E_z[X(z)]\E_z[X(z)]^\top
    =
    b_a^o(b_a^o)^\top.
\]
Thus
\[
    \lambda_{\max}(\mathsf B_a^o)\ge\|b_a^o\|^2.
\]
\end{proof}

Finally, the visible tangent direction yields a same-density perturbation with negative second-order fitting curvature at the antipodal saddle.
The proposition is an existence statement for a negative-curvature direction.
It does not by itself prove that the Hessian-guided perturbation selects this
particular same-density direction; that selection issue belongs to the
perturbation analysis rather than to the geometric verification.

\begin{proposition}[Dirac-supported transverse descent from unweighted transverse imbalance]
\label{lemma:descent-direction-antipodal-dirac-from-unweighted-imbalance}
\textbf{Assumptions used.}
\cref{assumption:on-sphere,assumption:odd-symmetric-teacher,assumption:quantitative-unweighted-two-sided-transverse-imbalance}.

Assume \(\mathcal O\ne\varnothing\).  Set \(a=a_\star\), so that the
perturbed Dirac pair is the fixed regularizer target.  Let
\[
    \nu_a^\pm
    \coloneq
    \frac12\delta_a+\frac12\delta_{-a},
    \qquad
    \mu^\dagger
    \coloneq
    \nu_a^\pm\otimes\nu_a^\pm.
\]
By                                                                      
\cref{lem:quantitative-control-weighted-transverse-moment},
\cref{lem:active-weighted-transverse-moment-implies-visibility},
\[
    \gamma_{\mathrm{esc}}
    \coloneq
    \lambda_{\max}(\mathsf B_a^o)
    \ge
    \kappa_{\min}(r_{\rm vis})^2\gamma_{\perp}^2
    >0.
\]
Let \(s_\star\in a^\perp\) be a unit eigenvector associated with
\(\gamma_{\mathrm{esc}}\).
Then there exists a same-density spherical
pushforward perturbation
\[
    \mu_h=\mu_{1,h}\otimes\mu_{2,h},
    \qquad
    \mu_{1,h}=\mu_{2,h},
    \qquad
    \mu_{i,0}=\nu_a^\pm,
\]
such that
\[
    \frac{\dd}{\dd h}\Big|_{h=0}F(\mu_h)=0
\]
and
\[
    F(\mu_h)
    =
    F(\mu^\dagger)
    -
    h^2\gamma_{\mathrm{esc}}
    +
    o(h^2).
\]
In particular,
\[
    \frac{\dd^2}{\dd h^2}\Big|_{h=0}F(\mu_h)
    =
    -2\gamma_{\mathrm{esc}}<0.
\]
Moreover, for the metric regularizer
\[
    F_{\mathrm{reg}}(\mu_1,\mu_2)
    =
    \frac{\lambdaReg}{p_W}
    Z_\pm(\mu_1,\mu_2)^{p_W/2},
    \qquad
    p_W=2+\delta_W,
\]
one has
\[
    F_{\mathrm{reg}}(\mu_h)=O(|h|^{2+\delta_W})=o(h^2).
\]
Thus the metric regularizer does not affect the quadratic descent term at
the antipodal Dirac saddle.
\end{proposition}

\begin{proof}
\noindent\textbf{Step 1: compute the moments at the antipodal pair.}
We first compute the active odd moments at the antipodal Dirac pair.  For
every odd \(n\in\mathcal O\),
\[
    T_n(\nu_a^\pm)
    =
    \frac12 a^{\otimes n}
    +
    \frac12(-a)^{\otimes n}
    =
    0.
\]
Therefore
\[
    \Theta(\mu^\dagger)=(0,0).
\]


\par\noindent\textbf{Step 2: construct the visible spherical perturbation.}

Recall that
\[
 a^\perp
 =\{s\in\mathbb R^{d_w}:\langle s,a\rangle_{\mathbb R^{d_w}}=0\}
 =T_a\Sphere^{d_w-1}=T_{-a}\Sphere^{d_w-1},
 \qquad
 \Pi_{a^\perp}=I-aa^\top,
\]
where the projection is the one used in
\cref{def:two-sided-transverse-teacher-moment-measure}.  Thus the visible
unit vector \(s_\star\in a^\perp\) from the statement is a common tangent
direction at both poles.  Set
\begin{equation}
\label{eq:mf-visible-tangent-field}
 v(w):=\frac12(I-ww^\top)s_\star,\qquad
 u_+:=v(a)=\frac{s_\star}{2},\qquad
 u_-:=v(-a)=\frac{s_\star}{2}.
\end{equation}
Indeed, \(w^\top v(w)=0\), so \(v(w)\in T_w\Sphere^{d_w-1}\).
Moreover, substituting \(w=a\) and \(w=-a\) and using
\(a^\top s_\star=0\) gives
\[
 v(a)=\tfrac12(I-aa^\top)s_\star=\tfrac12s_\star,
 \qquad
 v(-a)=\tfrac12(I-aa^\top)s_\star=\tfrac12s_\star.
\]
Hence the two atoms move in the same ambient transverse direction and
\(u_++u_-=s_\star\).  The factor \(1/2\) in \(v\) compensates for the
two atom masses \(1/2\), so that their summed first-order moment change is
exactly the differential \(\mathcal U_a(s_\star)\) defined in
\cref{def:tangential-visibility-matrix-dirac-pair}.
Use it in both marginals and define
\(\mu_{1,h}=\mu_{2,h}:=(R_h)_\#\nu_a^\pm\), where, for
\(w\in\{\pm a\}\),
\[
    R_h(w)=\frac{w+h v(w)}{\|w+h v(w)\|}.
\]
Since \(v(w)\in T_w\Sphere^{d_w-1}\), we have \(w^\top v(w)=0\), and hence
\[
    \frac{\dd}{\dd h}\Big|_{h=0}R_h(w)=v(w).
\]
Thus
\[
    \frac{\dd}{\dd h}\Big|_{h=0}R_h(w)^{\otimes n}
    =
    n\,\sym{w^{\otimes(n-1)}\otimes v(w)}.
\]
Consequently,
\begin{align}
    \dot T_n[v]
    &\coloneq
    \frac{\dd}{\dd h}\Big|_{h=0}T_n((R_h)_\#\nu_a^\pm)
    \\
    &=
    \frac n2
    \sym{a^{\otimes(n-1)}\otimes u_+}
    +
    \frac n2
    \sym{(-a)^{\otimes(n-1)}\otimes u_-}.
\end{align}
Because \(n\) is odd, \(n-1\) is even, and therefore
\[
    (-a)^{\otimes(n-1)}=a^{\otimes(n-1)}.
\]
Hence
\[
    \dot T_n[v]
    =
    \frac n2
    \sym{a^{\otimes(n-1)}\otimes(u_++u_-)}
    =
    \frac n2
    \sym{a^{\otimes(n-1)}\otimes s_\star}
    =
    \mathcal U_a(s_\star)_n.
\]
Since \(\mu_{1,h}=\mu_{2,h}\), we have
\[
    \dot{\mathbf S}[v]
    =
    \frac{\dot{\mathbf T}_1[v]+\dot{\mathbf T}_2[v]}{2}
    =
    \mathcal U_a(s_\star),
    \qquad
    \dot{\mathbf A}[v]
    =
    \frac{\dot{\mathbf T}_1[v]-\dot{\mathbf T}_2[v]}{2}
    =
    0.
\]
Thus this path leaves the active origin in the symmetric coordinate
\(\mathbf S\) with velocity \(\mathcal U_a(s_\star)\), while its
anti-symmetric coordinate \(\mathbf A\) remains zero.


\par\noindent\textbf{Step 3: expand the fitting term.}
Now we compute the model block.  By the active Hermite product formula,
\[
    \widetilde{\mathcal M}_p(\mathbf S,\mathbf A)
    =
    \sum_{\substack{
        n,m\in\mathcal O\\
        0\le r\le \min(n,m)\\
        p=n+m-2r
    }}
    c_{n,m,r}(a_na_m^\top)
    \otimes
    \sym{
        (S_n+A_n)\contr{r}(S_m-A_m)
    }.
\]
Along the same-density perturbation above,
\[
    \mathbf A(h)=0
\]
for every \(h\). The normalized atom maps constructed from the field in
\cref{eq:mf-visible-tangent-field} are smooth near $h=0$, so their finitely
many tensor moments are $C^2$ in $h$. Step 1 gives
$\mathbf S(0)=0$, and the derivative calculation in Step 2 gives
$\dot{\mathbf S}(0)=\mathcal U_a(s_\star)$. Taylor's formula therefore gives
\[
    \mathbf S(h)
    =
    h\,\mathcal U_a(s_\star)+O(h^2).
\]
Therefore
\begin{align}
    \widetilde{\mathcal M}_p(\Theta(\mu_h))
    &=
    \widetilde{\mathcal M}_p(\mathbf S(h),0)
    \\
    &=
    \sum_{\substack{
        n,m\in\mathcal O\\
        0\le r\le \min(n,m)\\
        p=n+m-2r
    }}
    c_{n,m,r}(a_na_m^\top)
    \otimes
    \sym{S_n(h)\contr{r}S_m(h)}
    \\
    &=
    h^2
    \underbrace{\sum_{\substack{
        n,m\in\mathcal O\\
        0\le r\le \min(n,m)\\
        p=n+m-2r
    }}
    c_{n,m,r}(a_na_m^\top)
    \otimes
    \sym{
        \mathcal U_a(s_\star)_n
        \contr{r}
        \mathcal U_a(s_\star)_m
    }}_{\eqqcolon  \mathcal Q_{p,+}(\mathcal U_a(s_\star))}
    +
    o(h^2).
\end{align}

We now expand the fitting term.  Recall that
\[
    F_{\mathrm{fit}}(\mu_h)
    =
    \frac12
    \sum_{p\ge0}
    \left\|
        \widetilde{\mathcal M}_p(\Theta(\mu_h))
        -
        \mathcal M_p^o
    \right\|_{\HS}^2.
\]
Using
\[
    \widetilde{\mathcal M}_p(\Theta(\mu_h))
    =
    h^2\mathcal Q_{p,+}(\mathcal U_a(s_\star))
    +
    o(h^2),
\]
we get
\begin{align}
    \frac12
    \left\|
        \widetilde{\mathcal M}_p(\Theta(\mu_h))
        -
        \mathcal M_p^o
    \right\|_{\HS}^2
    &=
    \frac12
    \left\|
        -\mathcal M_p^o
        +
        h^2\mathcal Q_{p,+}(\mathcal U_a(s_\star))
        +
        o(h^2)
    \right\|_{\HS}^2
    \\
    &=
    \frac12\|\mathcal M_p^o\|_{\HS}^2
    -
    h^2
    \left\langle
        \mathcal Q_{p,+}(\mathcal U_a(s_\star)),
        \mathcal M_p^o
    \right\rangle_{\HS}
    +
    o(h^2).
\end{align}
Summing over \(p\ge0\) proves
\[
    F_{\mathrm{fit}}(\mu_h)
    =
    F_{\mathrm{fit}}(\mu^\dagger)
    -
    h^2
    \sum_{p\ge0}
    \left\langle
        \mathcal Q_{p,+}(\mathcal U_a(s_\star)),
        \mathcal M_p^o
    \right\rangle_{\HS}
    +
    o(h^2).
\]


\par\noindent\textbf{Step 4: identify the visibility coefficient.}
It remains to identify the escape coefficient with the visibility matrix.
Use the first-order student feature $g_{a,s}(z)$ defined by the finite
Hermite sum in \eqref{eq:mf-visibility-first-order-feature}.
By the definition of the Hermite active blocks,
\[
    g_{a,s}(z)g_{a,s}(z)^\top
\]
has \(p\)-th active Hermite block
\[
    \mathcal Q_{p,+}(\mathcal U_a(s)).
\]
Similarly,
\[
    h^o(z)h^o(z)^\top
\]
has \(p\)-th active Hermite block \(\mathcal M_p^o\).  Therefore, by
orthogonality of the Hermite blocks,
\begin{align}
    \sum_{p\ge0}
    \left\langle
        \mathcal Q_{p,+}(\mathcal U_a(s)),
        \mathcal M_p^o
    \right\rangle_{\HS}
    &=
    \mathbb E_z
    \left[
        \left\langle
            g_{a,s}(z)g_{a,s}(z)^\top,
            h^o(z)h^o(z)^\top
        \right\rangle_{\mathrm F}
    \right]
    \\
    &=
    \mathbb E_z
    \left[
        \bigl(h^o(z)^\top g_{a,s}(z)\bigr)^2
    \right].
\end{align}
Now write \(g_{a,s}(z)=J_a(z)s\).  Then
\begin{align}
    \mathbb E_z
    \left[
        \bigl(h^o(z)^\top g_{a,s}(z)\bigr)^2
    \right]
    &=
    \mathbb E_z
    \left[
        s^\top J_a(z)^\top h^o(z)h^o(z)^\top J_a(z)s
    \right]
    \\
    &=
    s^\top
    \mathbb E_z
    \left[
        J_a(z)^\top h^o(z)h^o(z)^\top J_a(z)
    \right]
    s
    \\
    &=
    s^\top\mathsf B_a^o s.
\end{align}
Taking \(s=s_\star\) gives
\[
    \sum_{p\ge0}
    \left\langle
        \mathcal Q_{p,+}(\mathcal U_a(s_\star)),
        \mathcal M_p^o
    \right\rangle_{\HS}
    =
    s_\star^\top\mathsf B_a^o s_\star
    =
    \gamma_{\mathrm{esc}}.
\]
Hence
\[
    F_{\mathrm{fit}}(\mu_h)
    =
    F_{\mathrm{fit}}(\mu^\dagger)
    -
    h^2\gamma_{\mathrm{esc}}
    +
    o(h^2).
\]


\par\noindent\textbf{Step 5: compute the regularizer contribution.}

Recall that \(s_\star\) is the chosen unit transverse direction,
\(s_\star^\top a_\star=0\), and \(u_\pm=s_\star/2\). At the assigned
pole \(y=\pm a_\star\), the normalized perturbation satisfies
\[
 R_h(y)=\frac{y+(h/2)s_\star}{\sqrt{1+h^2/4}},\qquad
 y^\top R_h(y)=\frac1{\sqrt{1+h^2/4}}.
\]
The angle lies in \([0,\pi/2)\). Because
\(\cos(\arctan x)=1/\sqrt{1+x^2}\) for \(x\ge0\),
\[
 d_\Sphere(R_h(y),y)
 =\arccos\frac1{\sqrt{1+h^2/4}}=\arctan(|h|/2).
\]
For sufficiently small \(|h|\) both atoms retain their cap assignments.
Each marginal has two atoms of mass \(1/2\), and there are two marginals;
therefore \cref{eq:fixed-assignment-W2-square-formula} gives
\begin{equation}
\label{eq:mf-visible-path-metric-distance}
 Z_\pm(\mu_h)
 =\sum_{i=1}^2\sum_{\varepsilon\in\{-1,1\}}\frac12\arctan^2(|h|/2)
 =2\arctan^2(h/2)=O(h^2).
\end{equation}
For \(h\ne0\), the regularizer has the explicit factorization
\[
 F_{\mathrm{reg}}(\mu_h)=\frac{\lambdaReg}{p_W}|h|^{p_W}
      \left(\frac{2\arctan^2(h/2)}{h^2}\right)^{p_W/2}.
\]
The factor in parentheses extends to a smooth positive function at zero,
with value \(1/2\), because \(\arctan(h/2)=h/2+O(h^3)\).
The derivative calculation following \cref{eq:mf-zero-distance-retraction}
then gives \(F_{\mathrm{reg}}'(h)=O(|h|^{p_W-1})\) and
\(F_{\mathrm{reg}}''(h)=O(|h|^{p_W-2})\).
Since \(p_W=2+\delta_W>2\), the regularizer is \(C^2\) at zero, both
derivatives there are zero, and \(F_{\mathrm{reg}}(\mu_h)=o(h^2)\).
Thus the ordinary second derivative of the full objective exists and equals
\(-2\gamma_{\mathrm{esc}}\), as well as satisfying the displayed expansion.
Combining the fitting expansion and the regularizer estimate proves the claim.
\end{proof}
